% !TeX program = pdflatex
% !TeX root = documento.tex
% Revisão acadêmica — capítulos 2 e 3; modelo ufctex entregue pelo aluno.
\documentclass[a4paper,12pt,chapter=TITLE,section=Title,subsection=Title,oneside,english,spanish,brazil]{abntex2}
% Para utilizar este template siga o tutorial disponível em http://www.biblioteca.ufc.br/wp-content/uploads/2015/09/tutorial-sharelatex.pdf

%%%%%%%%%%%%%%%%%%%%%%%%%%%%%%%%%%%%%%%%%%%%%%%%%%%%%%%
%% Você deve criar uma conta no Overleaf. Depois,    %%
%% vá nas opções no canto esquerdo superior da tela  %%
%% e clique em "Copiar Projeto". Dê um novo nome pa- %%
%% ra o projeto.                                     %%
%%                                                   %%
%% Os principais desenvolvedores deste template são: %%
%%                                                   %%
%%            Ednardo Moreira Rodrigues              %%
%%       (Doutor em Engenharia Elétrica - UFC)       %%
%%                      &                            %%
%%            Alan Batista de Oliveira               %%
%%           (Engenheiro Eletricista - UFC)          %%
%%                                                   %%
%% Revisão:                                          %%
%%                                                   %%
%% - Francisco Edvander Pires Santos;                %%
%% - Juliana Soares Lima;                            %%
%% - Izabel Lima dos Santos;                         %%
%% - Kalline Yasmin Soares Feitosa.                  %%
%% - Eliene Maria Vieira de Moura;                   %%
%%                                                   %%
%% Colaboradores                                     %%
%%                                                   %%
%% -Andrei Bosco Bezerra Torres                      %% 
%% (Professor - Sistemas e Mídias Digitais -         %%
%% Instituto Universidade Virtual - UFC)             %%
%% Tiago ALves Lima                                  %% 
%% (Aluno de Mestrado em Eng. Elétrica)              %%
%%                                                   %%
%% Grande parte do trabalho foi adaptado do template %%
%% da UECE elaborado por:                            %%
%% Thiago Nascimento  (UECE)                         %%
%% Project available on:                             %%
%% https://github.com/thiagodnf/uecetex2             %%
%%                                                   %%
%% "Dúvidas, esclarecimentos ou sugestões podem ser  %%
%% enviadas para o seguinte e-mail:                  %%
%%                                                   %%
%%             atendimentobch@ufc.br                 %%
%%                                                   %%
%% As últimas atualizações estão descritas no inicio %%
%% do arquivo "README.md".                           %%
%%                                                   %%
%%%%%%%%%%%%%%%%%%%%%%%%%%%%%%%%%%%%%%%%%%%%%%%%%%%%%%%

% Importações de pacotes
\usepackage[utf8]{inputenc}                         % Acentuação direta
\usepackage[T1]{fontenc}                            % Codificação da fonte em 8 bits
\usepackage{graphicx}                               % Inserir figuras
\usepackage{amsfonts, amssymb, amsmath}             % Fonte e símbolos matemáticos
\usepackage{booktabs}                               % Comandos para tabelas
\usepackage{verbatim}                               % Texto é interpretado como escrito no documento
\usepackage{multirow, array}                        % Múltiplas linhas e colunas em tabelas
\usepackage{indentfirst}                            % Endenta o primeiro parágrafo de cada seção.
\usepackage{listings}                               % Utilizar codigo fonte no documento
\usepackage{xcolor}
\usepackage{microtype}                              % Para melhorias de justificação?
\usepackage[portuguese,ruled,lined]{algorithm2e}    % Escrever algoritmos
\usepackage{algorithmic}                            % Criar Algoritmos  
%\usepackage{float}                                 % Utilizado para criação de floats
\usepackage{amsgen}
\usepackage{lipsum}                                 % Usar a simulação de texto Lorem Ipsum
%\usepackage{titlesec}                              % Permite alterar os títulos do documento
\usepackage{tocloft}                                % Permite alterar a formatação do Sumário
\usepackage{etoolbox}                               % Usado para alterar a fonte da Section no Sumário
\usepackage[nogroupskip,nonumberlist]{glossaries}   % Permite fazer o glossario

\usepackage[font=singlespacing,format=plain,justification=justified,skip=0pt,singlelinecheck = false]{caption}            % Altera o comportamento da tag caption

\usepackage[alf, abnt-emphasize=bf, recuo=0cm, abnt-etal-cite=2, abnt-etal-list=0, abnt-etal-text=it]{abntex2cite}  % Citações padrão ABNT
%\usepackage[bottom]{footmisc}                      % Mantém as notas de rodapé sempre na mesma posição
%\usepackage{times}                                 % Usa a fonte Times
%%%%%%%%%%%%%%%%%%% AVISO %%%%%%%%%%%%%%%%%%%%%%%%%%%%%%%%%%%%%%%%
%descomente as duas linhas abaixo para alterar o texto de Times New Roman para Arial:

%\usepackage{helvet}
%\renewcommand{\familydefault}{\sfdefault}  % Usa a fonte Arial              
%%%%%%%%%%%%%%%%%%%%%%%%%%%%%%%%%%%%%%%%%%%%%%%%%%%%%%%%%%%%%%%%%%

\usepackage{mathptmx}         % Usa a fonte Times New Roman			%\usepackage{lmodern}         % Usa a fonte Latin Modern
%\usepackage{subfig}          % Posicionamento de figuras
%\usepackage{scalefnt}        % Permite redimensionar tamanho da fonte
%\usepackage{color, colortbl} % Comandos de cores
%\usepackage{lscape}          % Permite páginas em modo "paisagem"
%\usepackage{ae, aecompl}     % Fontes de alta qualidade
%\usepackage{picinpar}        % Dispor imagens em parágrafos
%\usepackage{latexsym}        % Símbolos matemáticos
%\usepackage{upgreek}         % Fonte letras gregas
\usepackage{appendix}         % Gerar o apendice no final do documento
\usepackage{paracol}          % Criar paragrafos sem identacao
\usepackage{lib/ufctex}	      % Biblioteca com as normas da UFC para trabalhos academicos
\usepackage{pdfpages}         % Incluir pdf no documento
\usepackage{amsmath}          % Usar equacoes matematicas

\makeglossaries % Organiza e gera a lista de abreviaturas, simbolos e glossario
\makeindex      % Gera o Indice do documento         



\trabalhoacademico{tccgraduacao}
\ehqualificacao{nao}
\hypersetup{hidelinks}
\ies{Universidade Federal do Ceará}
\iessigla{UFC}
\centro{Campus de Sobral}
\departamento{}
\graduacaoem{Engenharia Elétrica}
\habilitacao{bacharel}
\autor{Francisco de Assis Alves Vasconcelos Filho}
\titulo{Análise comparativa das especificações técnicas CNC-OMBR-MAT-18-0126-EDCE e CNC-NDBR-DBR-25-1580-EDBR aplicadas ao fornecimento de energia elétrica em empreendimentos com múltiplas unidades consumidoras}
\data{2026}
\local{Sobral}
\orientador{Prof. Eber de Castro Diniz}
\orientadories{Universidade Federal do Ceará (UFC)}
\orientadorcentro{Campus de Sobral}
\orientadorfeminino{nao}
\coorientador{}
\membrodabancadois{}
\membrodabancatres{}
\membrodabancaquatro{}
\membrodabancacinco{}
\membrodabancaseis{}
\begin{document}
\pretextual
\imprimircapa
\imprimirfolhaderosto{}
% Os elementos seguintes dependem da conclusão do TCC:
% ficha catalográfica, folha de aprovação, resumo, abstract e listas.
\imprimirsumario
\textual
\chapter{Introdução}
\label{cap:introducao}

\chapter{Objetivos}
\label{cap:objetivos}

\section{Objetivo geral}
\section{Objetivos específicos}
\section{Organização do trabalho}

\chapter{Trabalhos relacionados}
\label{cap:trabalhos-relacionados}

\chapter{Fundamentação teórica}
\label{cap:fundamentacao}

A conexão de empreendimentos com múltiplas unidades consumidoras ao sistema de distribuição requer a definição das condições de atendimento, do dimensionamento da entrada de energia e da organização da medição. Essas decisões são condicionadas pela regulamentação do setor elétrico e pelas especificações técnicas da distribuidora responsável pela área de concessão. No Ceará, os requisitos anteriormente estabelecidos pela especificação técnica CNC-OMBR-MAT-18-0126-EDCE foram substituídos pelos da especificação técnica CNC-NDBR-DBR-25-1580-EDBR, documento de abrangência nacional no âmbito da Enel. Neste trabalho, esses documentos serão denominados, respectivamente, especificação 0126 e especificação 1580. A avaliação dessa transição exige distinguir disposições regulatórias, critérios construtivos e requisitos regionais.

\section{Regulamentação do fornecimento de energia elétrica}
\label{sec:regulamentacao}

As condições gerais de prestação do serviço público de distribuição são regulamentadas pela Agência Nacional de Energia Elétrica (ANEEL). A Resolução Normativa ANEEL nº 1.000/2021 consolida disposições sobre conexão, atendimento e responsabilidades dos agentes envolvidos. Os critérios relativos ao grupo e à tensão de fornecimento devem, portanto, ser interpretados conjuntamente com as normas técnicas da distribuidora, sem afastar as disposições regulatórias aplicáveis. A Resolução Normativa ANEEL nº 1.059/2023 promoveu alterações em dispositivos da regulamentação original, conforme \citeonline{aneel1000} e \citeonline{aneel1059}.

O enquadramento das unidades consumidoras nos Grupos A e B não decorre exclusivamente da potência demandada pelo empreendimento. De acordo com o art. 23 da Resolução Normativa ANEEL nº 1.000/2021, a carga instalada, as características da rede de distribuição e as hipóteses regulamentares de atendimento constituem parâmetros para definição da tensão de conexão. A referência geral de 75~kW para determinadas conexões em tensão inferior a 2,3~kV não elimina as possibilidades excepcionais previstas na regulamentação. Por essa razão, limites atribuídos às unidades consumidoras individuais não devem ser confundidos com os limites de demanda da entrada coletiva.

A entrada de energia está sujeita, ainda, aos requisitos de segurança e desempenho das instalações elétricas. A ABNT NBR 5410 e outros documentos técnicos referenciados pelas especificações da Enel integram esse referencial, respeitadas as respectivas áreas de aplicação. A análise documental realizada pela distribuidora não transfere ao agente responsável pelo serviço público a responsabilidade técnica do profissional habilitado pelo projeto e pelas instalações particulares.

\section{Conexão coletiva e múltiplas unidades consumidoras}
\label{sec:conexaocoletiva}

A conexão coletiva caracteriza-se pelo atendimento a um conjunto de unidades consumidoras mediante infraestrutura de entrada e distribuição concebida para o empreendimento. Em edifícios residenciais, os circuitos destinados às unidades autônomas coexistem com circuitos associados aos serviços de uso comum. Essa configuração requer a compatibilização dos pontos de conexão, da distribuição interna, dos dispositivos de proteção e das medições individualizadas.

Na especificação 0126, o conceito de prédio de múltiplas unidades consumidoras (PMUC) orienta requisitos associados à concentração das medições e à alimentação dos consumidores vinculados à mesma edificação. A especificação abrange unidades autônomas e instalações das áreas comuns, estabelecendo condições próprias para a entrada de serviço, a apresentação do projeto e o arranjo dos equipamentos de proteção e medição, conforme \citeonline{enel0126}.

A especificação 1580 utiliza a denominação \emph{conexão coletiva} e abrange empreendimentos horizontais e verticais. O item 7.8.1 distingue dessa condição as edificações contíguas ou geminadas sem área de uso comum e alinhadas à via pública, às quais se aplicam critérios de conexão individual. A delimitação da modalidade de atendimento depende, desse modo, das características funcionais e elétricas do conjunto, não apenas de sua configuração arquitetônica, conforme \citeonline[item~7.8.1]{enel1580}.

\section{Carga instalada e demanda no dimensionamento da entrada}
\label{sec:demandas}

A carga instalada corresponde à soma das potências nominais dos equipamentos previstos ou existentes na instalação, de acordo com os critérios aplicáveis a cada categoria de consumidor. A demanda expressa a solicitação de potência em determinada condição de operação. Em empreendimentos com várias unidades consumidoras, a simultaneidade de utilização dos equipamentos constitui parâmetro essencial do dimensionamento da entrada geral, uma vez que a soma aritmética das cargas individuais não representa, necessariamente, a maior potência requerida ao sistema.

A especificação 0126 apresenta, em sua seção 7.11, critérios de estimativa de demanda associados à área útil das unidades residenciais e procedimentos destinados a outras categorias de carga. A aplicação desses critérios depende da configuração das unidades e das características das cargas previstas. Na especificação 1580, a determinação da demanda continua vinculada ao profissional habilitado, com referência às metodologias recomendadas pela distribuidora e aos documentos técnicos associados. Os dispositivos correspondentes estão apresentados na seção~7.11 da especificação 0126 e no item~7.3.2 da especificação 1580, conforme \citeonline{enel0126} e \citeonline{enel1580}.

A distinção entre carga instalada, demanda aparente e corrente de projeto é necessária para evitar equivalências indevidas entre grandezas. Potências expressas em quilowatts (kW) e quilovolt-ampères (kVA) não correspondem à mesma grandeza, assim como a corrente de projeto, em ampères (A), depende da potência aparente e do nível de tensão. Para uma alimentação trifásica equilibrada, a relação entre essas grandezas pode ser expressa por
\begin{equation}
 I = \frac{S}{\sqrt{3}\,V_{LL}},
 \label{eq:corrente-trifasica}
\end{equation}
na qual $I$ representa a corrente de linha, em ampères; $S$, a potência aparente trifásica, em volt-ampères; e $V_{LL}$, a tensão entre fases, em volts. Essa relação não substitui os procedimentos de dimensionamento de condutores e dispositivos de proteção, que também dependem das condições de instalação, da capacidade de condução de corrente, da queda de tensão e das correntes de falta previstas.

\section{Componentes da entrada coletiva e critérios técnicos}
\label{sec:componentes}

A entrada coletiva compreende instalações destinadas à conexão com a rede da distribuidora e à alimentação dos circuitos correspondentes às unidades consumidoras. Sua configuração pode incluir ramal de conexão, ramal de entrada, centro de proteção geral, barramentos, dispositivos de seccionamento, quadros de distribuição e centros de medição. A terminologia e os limites de responsabilidade devem ser observados conforme o documento técnico aplicável, pois a localização física de determinado componente não estabelece, isoladamente, sua propriedade ou a responsabilidade por sua instalação.

O centro de medição concentra os equipamentos utilizados na medição individualizada das unidades consumidoras e pode incorporar módulos de proteção, distribuição e passagem. Sua disposição física influencia as condições de acesso, manutenção, identificação de circuitos e leitura pela distribuidora. A especificação 1580 dedica seções próprias à medição e aos centros de medição, acompanhadas de tabelas e desenhos de arranjos construtivos, conforme \citeonline[seções~7.17 a~7.21]{enel1580}.

A seleção dos dispositivos de proteção e a definição do sistema de aterramento integram o dimensionamento da entrada. A proteção geral deve ser compatível com a corrente de projeto, com a capacidade de interrupção requerida e com as condições de falta. As medidas de aterramento e equipotencialização dependem, por sua vez, das características da instalação e do esquema de alimentação. A especificação 1580 reúne prescrições sobre esses temas nas seções 7.24 e 7.25, sem afastar a necessidade de observância das normas brasileiras correlatas, conforme \citeonline{enel1580}.

\section{Especificações técnicas da Enel utilizadas na pesquisa}
\label{sec:especificacoes}

\subsection{Especificação CNC-OMBR-MAT-18-0126-EDCE}
\label{subsec:0126}

A especificação 0126, intitulada \emph{Fornecimento de energia elétrica a prédios de múltiplas unidades consumidoras}, foi utilizada pela Enel Distribuição Ceará para estabelecer critérios de projeto e execução das entradas de energia em PMUC. A versão 03, de 31 de março de 2022, é adotada como referência documental anterior nesta pesquisa. O documento abrange condições de fornecimento, medição, proteção, aterramento, cálculo de demanda e documentação necessária à análise dos projetos, conforme \citeonline{enel0126}.

A especificação inclui disposições aplicáveis a situações particulares e critérios condicionados ao local de atendimento. Sua interpretação requer a leitura conjunta do texto, dos desenhos técnicos vinculados ao documento e dos comunicados técnicos cuja aplicação possa ser demonstrada. A existência de disposições regulatórias anteriores a 2025 impede que requisitos já incorporados à versão de 2022 sejam classificados automaticamente como inovações introduzidas pela especificação nacional.

\subsection{Especificação CNC-NDBR-DBR-25-1580-EDBR}
\label{subsec:1580}

A especificação 1580, intitulada \emph{Fornecimento de energia elétrica em tensão secundária de distribuição até 1 kV -- conexão coletiva}, estabelece requisitos técnicos para o atendimento coletivo no âmbito das distribuidoras da Enel. A versão 2, de 8 de setembro de 2025, registra a substituição da especificação 0126, entre outros documentos regionais sobre o mesmo tema, conforme \citeonline[item~2]{enel1580}.

Apesar de sua abrangência nacional, a especificação contém condições particulares para as áreas de concessão do Ceará, do Rio de Janeiro e de São Paulo. Sua aplicação exige, portanto, identificar as prescrições compatíveis com o sistema de distribuição e com as soluções admitidas na região estudada. A reunião de requisitos regionais em um único documento não implica uniformidade das tensões nominais, modalidades de conexão ou alternativas construtivas em todas as áreas de concessão.

\section{Vigência e transição das especificações técnicas}
\label{sec:transicao}

A vigência de normas técnicas da distribuidora não se confunde necessariamente com a data registrada em sua folha de identificação. O art. 20 da Resolução Normativa ANEEL nº 1.000/2021 estabelece procedimentos para comunicação e entrada em vigor de novas normas ou de suas alterações. Como regra geral, o §~2º prevê prazo mínimo de 120 dias contado da comunicação aos interessados, ressalvadas as hipóteses excepcionais nele previstas. O §~4º também contempla a preservação de situações constituídas antes da vigência do novo padrão, nas condições regulamentares, conforme \citeonline[art.~20]{aneel1000}.

A data de aprovação de um projeto não identifica, por si só, todas as prescrições aplicadas durante sua análise. Para os fins desta pesquisa, a especificação 0126 constitui a referência normativa anterior, e a especificação 1580 constitui o documento de comparação posterior. A aplicação prática terá como referência a documentação de um empreendimento residencial de múltiplas unidades consumidoras, sem o propósito de reconstituir o histórico administrativo de aprovação ou verificar a execução das instalações.

\chapter{Metodologia}
\label{cap:metodologia}

A pesquisa adota procedimentos de análise documental e comparação técnica dos requisitos de conexão coletiva, complementados por aplicação a um projeto elétrico aprovado pela distribuidora. A unidade de análise normativa corresponde às exigências relativas a empreendimentos com múltiplas unidades consumidoras na área de concessão da Enel Ceará. O estudo de caso permite examinar as implicações das diferenças identificadas sobre soluções de engenharia documentadas, sem constituir avaliação de conformidade integral do empreendimento.

\section{Natureza e abordagem da pesquisa}
\label{sec:tipo-pesquisa}

Trata-se de pesquisa aplicada, descritiva e comparativa, fundamentada em documentos técnicos. A abordagem é predominantemente qualitativa, uma vez que o objetivo central consiste em identificar, classificar e interpretar mudanças dos requisitos normativos. Dados quantitativos são empregados nas situações em que a avaliação de repercussões sobre o projeto exige a confrontação de limites de potência, demanda, corrente, dimensões ou características nominais dos equipamentos.

A comparação das especificações e o estudo de caso constituem etapas relacionadas, porém distintas. Na primeira, identificam-se as diferenças entre os documentos de referência; na segunda, examina-se a aplicabilidade dessas diferenças às características de uma instalação projetada. A ausência de determinado requisito em uma seção não será interpretada, isoladamente, como sua eliminação, sendo necessária a verificação de remissões a outras seções, tabelas, anexos e documentos complementares.

\section{Conjunto documental}
\label{sec:documentos}

As fontes normativas principais são a especificação 0126, versão 03, de 31 de março de 2022, e a especificação 1580, versão 2, de 8 de setembro de 2025. Os documentos complementares serão consultados apenas quando necessários à interpretação dos requisitos selecionados.

A Resolução Normativa ANEEL nº 1.000/2021 e suas alterações pertinentes compõem o referencial regulatório para interpretação das condições de atendimento, das normas técnicas da distribuidora e dos períodos de transição. Também integram o conjunto documental os desenhos e anexos das especificações quando contiverem exigências que complementem o texto principal.

A aplicação prática utiliza a documentação do projeto elétrico do Edifício Palatium, empreendimento residencial localizado em Sobral, Ceará. O conjunto analisado abrange o memorial de cálculo elétrico, revisão 04, as pranchas do sistema de entrada e alimentação, o registro de responsabilidade técnica e a carta de aprovação expedida pela Enel Distribuição Ceará em 25 de outubro de 2025. Os dados técnicos do empreendimento serão apresentados na análise do estudo de caso, e não nesta caracterização metodológica. O histórico das revisões anteriores do projeto não integra o objeto da pesquisa.

\section{Procedimentos de comparação normativa}
\label{sec:procedimento}

O levantamento será realizado em etapas sucessivas. Inicialmente, serão identificados objeto, área de aplicação e elementos técnicos equivalentes nas duas especificações. Em seguida, as prescrições serão agrupadas conforme os temas: enquadramento e limites de atendimento, cálculo de demanda, formas de conexão, ramais, padrão de entrada, medição, proteção, aterramento e exigências de apresentação de projeto.

Para cada tema, serão registrados os dispositivos correspondentes em cada documento, com indicação da seção, item, tabela ou desenho, quando aplicável. A comparação será efetuada entre prescrições funcionalmente equivalentes, inclusive quando inseridas em partes distintas das especificações. Quando o documento de 2025 remeter a outra norma, serão examinadas a referência e suas condições de aplicação antes de caracterizar eventual supressão ou modificação de requisito.

As diferenças serão classificadas em seis categorias: requisito mantido, alterado, acrescentado, retirado ou substituído, reorganização textual sem efeito técnico demonstrado e situação não conclusiva. A classificação estará vinculada às evidências documentais. A existência de nota regional, condição excepcional ou exigência de estudo de viabilidade será registrada expressamente, de modo a evitar generalizações que não se apliquem ao Ceará.

\section{Critérios de avaliação dos impactos técnicos}
\label{sec:impactos}

Os impactos serão avaliados segundo a possibilidade de interferência em decisões de engenharia documentadas. Serão considerados a modalidade de alimentação, a localização e a configuração do padrão de entrada, o dimensionamento de condutores e dispositivos de proteção, a disposição dos centros de medição e os requisitos de documentação e aprovação.

A classificação de determinada exigência como alterada não implicará, por si só, a necessidade de reformular o projeto utilizado como referência. Para estabelecer uma consequência técnica, será necessário identificar a solução registrada na documentação, o requisito correspondente na especificação 1580 e as condições de aplicação. Os resultados serão classificados como solução compatível com o requisito examinado, adequação potencialmente necessária, situação condicionada à avaliação da distribuidora ou informação documental insuficiente.

Quando disponíveis os parâmetros necessários, serão realizados cálculos demonstrativos com identificação das grandezas, hipóteses e fontes empregadas. Esses cálculos terão finalidade analítica e não substituirão os memoriais nem a responsabilidade técnica dos profissionais responsáveis pelo projeto original.

\section{Aplicação ao estudo de caso}
\label{sec:estudocaso}

A aplicação prática será conduzida sobre as soluções de entrada coletiva representadas no projeto aprovado do Edifício Palatium. A documentação existente será adotada como configuração de referência para verificar os possíveis efeitos das exigências selecionadas da especificação 1580 sobre um mesmo empreendimento. A especificação 0126 será considerada como referência normativa anterior, em conformidade com a delimitação da pesquisa.

A avaliação ficará restrita aos elementos presentes no memorial e nas pranchas do projeto. Para cada tema identificado como relevante na comparação normativa, será estabelecida uma relação entre a solução projetada, o requisito correspondente e a consequência técnica potencial. A pesquisa não promoverá alterações nos documentos originais nem atribuirá desconformidade ao projeto sem a confirmação das condições normativas pertinentes.

A carta de aprovação da concessionária comprova o deferimento do projeto apresentado, mas não discrimina os dispositivos normativos empregados na análise. O contexto de transição será examinado à luz das datas documentadas e do art. 20 da Resolução Normativa ANEEL nº 1.000/2021. Não serão presumidas a data efetiva de comunicação da especificação 1580 ou a incidência de hipóteses excepcionais sem suporte documental.

\section{Tratamento dos resultados e limitações}
\label{sec:limitacoes}

Os resultados da comparação serão sistematizados em quadros contendo os temas examinados, os dispositivos correspondentes das especificações 0126 e 1580, a natureza da alteração e a repercussão técnica identificada. Os resultados do estudo de caso serão apresentados separadamente, com indicação da solução documentada e das condições necessárias para cada avaliação.

A pesquisa limita-se às versões documentais selecionadas e aos critérios aplicáveis à área de concessão do Ceará. Não serão investigados os requisitos exclusivos de outras distribuidoras, o histórico de revisões internas do projeto, as instalações efetivamente executadas, os procedimentos de vistoria ou a energização do empreendimento. A análise não constitui laudo de conformidade nem substitui a avaliação técnica da concessionária em situações sujeitas a estudos de viabilidade.

Quando não houver elementos suficientes para determinada conclusão, a limitação será registrada no resultado correspondente. Esse procedimento preserva a distinção entre as informações expressamente documentadas, a interpretação de requisitos técnicos e as hipóteses dependentes de confirmação adicional.

\chapter{Resultados}
\label{cap:resultados}

\section{Análise comparativa das especificações 0126 e 1580}
\section{Estudo de caso: Edifício Palatium}

\chapter{Conclusões e trabalhos futuros}
\label{cap:conclusoes}

\postextual
\bibliography{3-pos-textuais/referencias}
\end{document}
